\documentclass[11pt,a4paper]{article}
\usepackage[utf8]{inputenc}
\usepackage[T1]{fontenc}
\usepackage[a4paper,margin=2.2cm]{geometry}
\usepackage{hyperref}
\usepackage{enumitem}
\usepackage{xcolor}
\hypersetup{colorlinks=true,linkcolor=blue,urlcolor=blue}
\title{MedCareSO\\Guide d'installation et de licence sécurisée}
\author{Procédure développeur et cabinet client}
\date{\today}
\begin{document}
\maketitle

\section*{But et limites}
MedCareSO est une application hors ligne. Aucune protection ne rend un logiciel
local impossible à modifier, mais cette procédure rend la copie de licence et
la falsification beaucoup plus coûteuses. La sécurité principale est une
licence JSON signée Ed25519 : la clé privée reste uniquement chez MedCareSO et
l'application client contient seulement la clé publique.

\section{Préparation unique chez MedCareSO}
\begin{enumerate}[label=\arabic*.]
  \item Sur un poste développeur sûr, créez une paire de clés :
\begin{verbatim}
npm run license:generate-keypair -- --out-dir /chemin/sur/disque-chiffre/medcareso
\end{verbatim}
  \item Conservez \texttt{medcareso-license-private.pem} dans un coffre chiffré.
  Ne l'envoyez jamais à un client, ne le placez jamais dans le dépôt Git, dans
  \texttt{assets}, ou dans le dossier de l'application.
  \item Installez uniquement la clé publique dans le code de la version à
  distribuer :
\begin{verbatim}
npm run license:install-public-key -- --key /chemin/sur/disque-chiffre/medcareso/medcareso-license-public.pem
\end{verbatim}
  \item Créez le paquet client signé. La variable pointe vers la clé privée
  stockée hors du dépôt :
\begin{verbatim}
export MEDCARESO_LICENSE_PRIVATE_KEY_FILE=/chemin/sur/disque-chiffre/medcareso/medcareso-license-private.pem
npm run build:secure
\end{verbatim}
Cette commande crée un manifeste d'intégrité signé puis le paquet Electron
ASAR. Une version client ne doit pas être distribuée avec \texttt{npm run build}
simple.
\end{enumerate}

\section{Installation chez un nouveau cabinet}
\begin{enumerate}[label=\arabic*.]
  \item Installez la version officielle construite avec \texttt{build:secure}.
  \item Démarrez MedCareSO. La fenêtre d'activation affiche l'\emph{Identifiant
  de ce poste} (empreinte SHA-256). Copiez exactement cette valeur et
  communiquez-la au développeur.
  \item Ne donnez pas de clé texte au cabinet. Le cabinet doit recevoir un
  fichier \texttt{.medcareso.json} signé.
  \item Chez MedCareSO, signez le fichier avec l'empreinte reçue. Exemple pour
  un abonnement se terminant le 18 juillet 2027 :
\begin{verbatim}
npm run license:sign -- \
  --private-key /chemin/sur/disque-chiffre/medcareso/medcareso-license-private.pem \
  --machine-id EMPREINTE_RECUE \
  --cabinet "Cabinet Dr. Exemple" \
  --expires 2027-07-18 \
  --features standard \
  --out /chemin/partage/Cabinet-Dr-Exemple.medcareso.json
\end{verbatim}
Pour une licence sans expiration, retirez \texttt{--expires}.
  \item Transmettez le fichier au cabinet par un canal convenu (clé USB,
  messagerie chiffrée ou portail client). Le fichier ne contient pas la clé
  privée.
  \item Dans MedCareSO, cliquez sur \emph{Choisir le fichier}, sélectionnez la
  licence, puis \emph{Activer la licence}. La signature, l'empreinte de la
  machine et l'expiration sont vérifiées avant l'enregistrement.
\end{enumerate}

\section{Renouvellement, changement de PC et incident}
\begin{itemize}
  \item Pour migrer une installation existante, connectez-vous une dernière fois
  avec le super administrateur, relevez l'identifiant du poste, puis installez
  le nouveau fichier signé. Les anciens codes texte ne sont plus acceptés.
  \item Pour renouveler un abonnement, demandez à nouveau l'identifiant du
  poste et signez un nouveau fichier avec une nouvelle date d'expiration.
  \item Pour un changement de PC, l'empreinte change. Créez donc un nouveau
  fichier pour le nouveau poste. Désactivez l'ancien fichier dans l'application
  si le PC est encore disponible.
  \item Une réinstallation de MedCareSO ne réinitialise pas la licence : elle
  est stockée dans les données utilisateur du système, hors du dossier
  d'installation.
  \item Si l'horloge système est remontée avant la dernière utilisation connue,
  l'application demande une intervention. Corrigez l'heure du système et
  contactez MedCareSO si le problème persiste.
\end{itemize}

\section{Protections présentes dans cette version}
\begin{itemize}
  \item \textbf{Electron :} isolation de contexte, Node désactivé dans le
  renderer, sandbox activée pour toutes les fenêtres, DevTools et raccourcis
  d'inspection désactivés dans les versions empaquetées.
  \item \textbf{Licence :} signature Ed25519 vérifiée dans le process principal,
  empreinte machine, vérification au démarrage et toutes les dix minutes,
  détection simple de retour arrière de l'horloge.
  \item \textbf{Intégrité :} les fichiers principaux sont hachés SHA-256 lors du
  build et le manifeste est signé. Une modification détectée bloque le
  démarrage de la version empaquetée.
  \item \textbf{Données :} cette version exploite PostgreSQL, pas SQLite ;
  SQLCipher ne s'applique donc pas. Utilisez un utilisateur PostgreSQL dédié,
  un mot de passe long, un pare-feu limitant PostgreSQL au réseau du cabinet,
  TLS si la base traverse un réseau, et des sauvegardes chiffrées.
\end{itemize}

\section{Obfuscation et exploitation}
La commande \texttt{build:secure} obfusque après l'empaquetage uniquement les
modules de licence et d'intégrité dans \texttt{app.asar}; le dépôt reste donc
lisible et aucun code obfusqué n'est commité. Ne stockez ni mot de passe de
base, ni clé privée, ni secret API dans le renderer. Testez toujours
l'installateur sur un PC vierge avant livraison.

\section*{Checklist avant remise client}
\begin{enumerate}[label={[\ ]}]
  \item Clé privée uniquement dans un coffre développeur.
  \item Clé publique installée dans le build.
  \item \texttt{MEDCARESO\_LICENSE\_PRIVATE\_KEY\_FILE} défini pour \texttt{build:secure}.
  \item Installateur généré, signature de licence vérifiée sur le poste client.
  \item Sauvegarde PostgreSQL chiffrée et accès réseau restreint.
  \item Identifiants administrateur par défaut modifiés.
\end{enumerate}
\end{document}
